\documentclass[../../../include/open-logic-section]{subfiles}

\begin{document}

\olfileid{sfr}{infinite}{card-sb}

\olsection{అనుబంధం: ష్రోడర్--బెర్న్‌స్టైన్ నిరూపణ}

సరళ సమితి సిద్ధాంతాన్ని విడిచిపెట్టేముందు, పరిశీలించాల్సిన చివరి
సరళమైన (కానీ అధునాతనమైన!) నిరూపణ ఒకటి ఉంది. ఇది
ష్రోడర్--బెర్న్‌స్టైన్ నిరూపణ
(\olref[sfr][siz][sb]{thm:schroder-bernstein}):
$\cardle{A}{B}$ మరియు $\cardle{B}{A}$ అయితే $\cardeq{A}{B}$;
అంటే, \tetoken{ఒకటి-ఒకటి ప్రమేయాలు}{!!{injection}s} $f \colon A \to B$, $g \colon B \to A$
ఇచ్చినప్పుడు \tetoken{ఒక ద్విగుణ ప్రమేయం}{!!a{bijection}} $h \colon A \to B$ ఉంటుంది.

ఈ అధ్యాయంలో డెడెకిండ్ చెప్పిన \emph{సంవృతాల} భావనను అనుసరించాం.
నిజానికి ఈ భావనను ఉపయోగించి డెడెకిండ్ ష్రోడర్--బెర్న్‌స్టైన్‌కు ఒక
అందమైన నిరూపణ ఇచ్చాడు; దానినే ఇక్కడ అందిస్తాం. కొద్దిగా భిన్నమైన,
కానీ సారాంశంగా సారూప్యమైన వివరణ కావాలంటే, ఈ నిరూపణ
\citet[pp.~157--8]{Potter2004}ను దగ్గరగా అనుసరిస్తుంది. కొద్దిగా
గూగుల్‌లో వెతికినా, $\sqrt{2}$ అకరణీయత మాదిరిగానే ఈ సిద్ధాంతానికి
\emph{అనేక} ఆసక్తికరమైన, భిన్నమైన నిరూపణలు ఉన్నాయని తెలుస్తుంది.

\olref[sfr][infinite][dedekind]{Closure}లోని సంజ్ఞకు సారూప్యంగా,
ఏ సమితి $U$, ప్రమేయం $f \colon U \to U$, ఉపసమితి $B \subseteq U$కైనా
ఇలా తీసుకుందాం:
\[
	\Closureofunder{f}{B} = \bigcap \Setabs{X \subseteq U}{B \subseteq X
	\text{ and $X$ is $f$-closed}}
\]
ఇలా నిర్వచించిన $\Closureofunder{f}{B}$, $B$ను కలిగిన అతి చిన్న
$f$-సంవృత సమితి; అంటే:
\sourcecorrection{OLTEINF-004}{ఆంగ్ల మూలం ఈ సమితి-సంవృతాన్ని ఏ
ప్రమేయం, ఏ సమితికైనా ప్రకటించి, $f(x)$ నిర్వచితమయ్యే పరిసర సమితిని
బంధించలేదు. తరువాతి వినియోగం నిర్ణయించినట్లుగా $f \colon U\to U$,
$B\subseteq U$, అభ్యర్థి $X\subseteq U$ అని రకాలను స్పష్టం చేశాం.}

\begin{lem}\ollabel{Closureprops}
	ఏ ప్రమేయం $f \colon U \to U$, ఏ $B \subseteq U$కైనా:
	\begin{enumerate}
		\item\ollabel{Closurehaselem} $B \subseteq \Closureofunder{f}{B}$; మరియు
		\item\ollabel{Closureclosed} $\Closureofunder{f}{B}$ $f$-సంవృతం; మరియు
		\item\ollabel{Closuresmallest} $X$ $f$-సంవృతమై, $B
		\subseteq X$ అయితే, $\Closureofunder{f}{B} \subseteq X$.
	\end{enumerate}
\end{lem}

\begin{proof}
\olref[sfr][infinite][dedekind]{closureproperties}లో చేసినట్లే.
\end{proof}

ష్రోడర్--బెర్న్‌స్టైన్‌ను చేరుకోవడానికి చివరి వాస్తవం ఒకటి కావాలి:

\begin{prop}\ollabel{sbhelper}
$A \subseteq B \subseteq C$ మరియు $A \approx C$ అయితే,
$\cardeq{B}{C}$.
\end{prop}
\sourcecorrection{OLTEINF-005}{ఆంగ్ల మూలపు ముగింపులో
$\cardeq{\cardeq{A}{B}}{C}$ అనే వికృతంగా గూడుకట్టిన పరిమాణ-సమానత
ఉంది. ఉపపత్తులు, తరువాతి $C$ నుంచి $B$కు ద్విగుణ ప్రమేయ నిరూపణ
నిర్ణయించినట్లుగా ముగింపును $\cardeq{B}{C}$గా సరిచేశాం.}

\begin{proof}
\tetoken{ఒక ద్విగుణ ప్రమేయం}{!!a{bijection}} $f \colon C \to A$ ఇచ్చినప్పుడు,
$F = \Closureofunder{f}{C \setminus B}$గా తీసుకుని, ప్రవేశం $C$ గల
$g$ అనే ప్రమేయాన్ని ఇలా నిర్వచించండి:
\[
	g(x) =
	\begin{cases}
			f(x) &\text{if $x \in F$}\\
			x & \text{otherwise}
		\end{cases}
\]
$g$, $C \to B$ నుంచి \tetoken{ఒక ద్విగుణ ప్రమేయం}{!!a{bijection}} అని చూపుతాం; దాని నుంచి
$\comp{f^{-1}}{g} \colon A \to B$ ఒక \tetoken{ద్విగుణ ప్రమేయం}{!!a{bijection}} అని వస్తుంది;
దీనితో నిరూపణ పూర్తవుతుంది.

మొదట $x \in F$, కానీ $y\notin F$ అయితే $g(x) \neq g(y)$ అని
వాదిస్తాం. వైరుధ్యం కోసం అలా కాదనుకుందాం; అప్పుడు
$y = g(y) = g(x) = f(x)$. $x \in F$, అలాగే
\olref{Closureprops} ప్రకారం $F$ $f$-సంవృతం కనుక
$y = f(x) \in F$; ఇది వైరుధ్యం.

ఇప్పుడు $g(x) = g(y)$ అనుకుందాం. పై వాదన ప్రకారం $x \in F$ అవడం
అంటే $y \in F$ అవడమే. $x, y \in F$ అయితే,
$f(x) = g (x) = g(y) = f(y)$; $f$ \tetoken{ఒక ద్విగుణ ప్రమేయం}{!!a{bijection}} కనుక $x = y$.
$x, y \notin F$ అయితే $x = g(x) = g(y) = y$. కాబట్టి $g$ ఒక
\tetoken{ఒకటి-ఒకటి ప్రమేయం}{!!a{injection}}.

ఇక $\ran{g} = B$ అని చూపాలి. మొదట ప్రతి $g$ విలువ $B$లోనే ఉంటుంది:
$y\in F$ అయితే $g(y)=f(y)\in A\subseteq B$; $y\notin F$ అయితే
$C\setminus B\subseteq F$ కనుక $y\in B$, అందువల్ల $g(y)=y\in B$.
తిరుగుగా $x \in B \subseteq C$ను స్థిరపరచండి. $x \notin F$ అయితే
$g(x) = x$. $x \in F$ అయితే, ఏదో $y \in F$కు $x = f(y)$;
లేకపోతే $F \setminus \{x\}$, $f$-సంవృతమై $C\setminus B$ను కలిగి
ఉండేది; ఇది \olref{Closureprops} ప్రకారం అసాధ్యం. ఇప్పుడు
$g(y) = f(y) = x$.
\sourcecorrection{OLTEINF-006}{ఆంగ్ల మూలం $\ran{g}=B$ అని చూపుతామని
చెప్పి, ప్రతి $x\in B$ను $g$ తాకుతుందని మాత్రమే చూపింది. ఇక్కడ
ముందుగా తప్పిపోయిన $\ran{g}\subseteq B$ దిశను, $F$లోని ఇన్‌పుట్లు
$A\subseteq B$లోకి వెళ్తాయని, $F$ బయటివి $C\setminus B\subseteq F$
వల్ల తప్పక $B$లోనే ఉంటాయని చూపి చేర్చాం.}
\end{proof}

చివరగా, ప్రధాన ఫలితపు నిరూపణ ఇదిగో. ఒక ప్రమేయం $h$, సమితి $D$
ఇచ్చినప్పుడు $\funimage{h}{D} = \Setabs{h(x)}{x \in D}$గా
నిర్వచిస్తామని గుర్తు చేసుకోండి.

\begin{proof}[ష్రోడర్--బెర్న్‌స్టైన్ నిరూపణ]
$f \colon A \to B$, $g \colon B \to A$లు \tetoken{ఒకటి-ఒకటి ప్రమేయాలు}{!!{injection}s} అనుకుందాం.
$\funimage{f}{A} \subseteq B$ కనుక
$\funimage{g}{\funimage{f}{A}} \subseteq g[B] \subseteq A$.
అంతేకాక $g$, $f$ రెండూ ఒకటి-ఒకటి కనుక
$\comp{f}{g} \colon A \to \funimage{g}{\funimage{f}{A}}$ ఒక
\tetoken{ఒకటి-ఒకటి ప్రమేయం}{!!{injection}}; నిజానికి దాని సహప్రవేశాన్ని నిర్వచించిన విధానం వల్లనే
$\comp{f}{g}$ ఒక \tetoken{ద్విగుణ ప్రమేయం}{!!a{bijection}}. అందువల్ల
$\cardeq{\funimage{g}{\funimage{f}{A}}}{A}$; కాబట్టి
\olref{sbhelper} ప్రకారం \tetoken{ఒక ద్విగుణ ప్రమేయం}{!!a{bijection}}
$h \colon A \to \funimage{g}{B}$ ఉంది. అంతేకాక $g^{-1}$ అనేది
$\funimage{g}{B} \to B$ నుంచి \tetoken{ఒక ద్విగుణ ప్రమేయం}{!!a{bijection}}. కాబట్టి
$\comp{h}{g^{-1}} \colon A \to B$ ఒక \tetoken{ద్విగుణ ప్రమేయం}{!!a{bijection}}.
\end{proof}

\end{document}
